\section{The question and the model}
\begin{frame}[t,label=question]{The question}
\small
Suppose a broker shows a client only some of their feasible options.
A better match may never reach them.

\medskip
We ask which firms a platform should audit to make its matching stable,
including against the options that were initially hidden.

\medskip
The difficulty is that opening an option can change the matching
algorithm's choice.
\sources{We model access to opportunities. Commissions and brokers' reasons for hiding options
are outside the model.}
\end{frame}

\begin{frame}[t,label=clients]{The market}
\small
Buyers $B$ and sellers $S$, with one property per seller.
Each client can have at most one partner.
\[
 E\subseteq B\times S
 \qquad\text{is the set of mutually acceptable, feasible pairs.}
\]
Each client strictly ranks its feasible partners. Every acceptable partner
is better than being unmatched, written $\bot$.

\medskip
Each client $i$ has one fixed representative $f(i)$.
A firm can represent many clients. A \emph{side-exclusive} firm represents
only buyers or only sellers.
\sources{Prices, rankings, and representation stay fixed. Brokers control access.
For rentals, interpret the two sides as renters and listing slots.}
\end{frame}

\begin{frame}[t,label=blocking]{Stability in the full market}
\small
For a matching $M$, let $M(i)$ be client $i$'s partner, or $\bot$ if unmatched.
\[
 \BP_E(M)=\{bs\in E\setminus M:
 s\succ_b M(b)\ \text{and}\ b\succ_s M(s)\}.
\]
A pair blocks when both clients would rather match with each other.

\medskip
We call $M$ \emph{full-market stable} when $\BP_E(M)=\varnothing$.

\medskip
Deferred acceptance is stable on the graph it receives.
A pair left out of that graph can still block in $E$.
\sources{This treatment of unavailable pairs follows restricted-edge matching
(Cseh and Manlove, 2016).}
\end{frame}

\begin{frame}[t,label=audit]{Audits and deferred acceptance}
\small
Initially, only $E_0\subseteq E$ is available. Let $\FF$ be the firms.
Auditing a set $A\subseteq\FF$ opens
\[
 E(A)=E_0\cup\{bs\in E:f(b)\in A\ \text{or}\ f(s)\in A\}.
\]
Either representative is enough. An audit enables all feasible pairs
involving its clients, without forcing a transaction.
\[
 c(A)=\sum_{f\in A}c_f,\quad c_f\in\mathbb Z_{>0},
 \qquad D(A)=\DA(E(A)).
\]
\textbf{Deferred acceptance (DA):} buyers propose down their available lists.
Each seller keeps its favorite proposal so far and rejects the others.
Rejected buyers continue until no further proposal is possible.
\sources{Assumption: the planner knows $E$ and the rankings. Payment enables access.
The platform always uses buyer-proposing DA.}
\end{frame}

\begin{frame}[t,label=objectives]{Two audit costs}
\small
\textbf{Response:} DA must select a full-market stable matching.
\[
 \resp=\min_{A\subseteq\FF}\{c(A):\BP_E(D(A))=\varnothing\}.
\]
\textbf{Support:} some full-market stable matching must be available.
\[
 \supp=\min_{A\subseteq\FF}
 \{c(A):\exists M\subseteq E(A),\ \BP_E(M)=\varnothing\}.
\]
Every successful DA response supplies such a matching, so
\[
 \supp\leq\resp.
\]
\sources{Auditing every firm is feasible. Running DA is polynomial-time.
Choosing the audits is the optimization problem.}
\end{frame}
